\vspace{0.5em}\noindent\textbf{Sharing and Feedback}\par

\vspace{0.5em}\noindent\textbf{Supervisor information}\par

{
\begingroup
\small
\setlength{\tabcolsep}{3pt}
\renewcommand{\arraystretch}{1.15}
\begin{longtable}{@{}ll@{}}
\toprule\noalign{}
Field & Value \\
\midrule\noalign{}
\endhead
\bottomrule\noalign{}
\endlastfoot
Supervisor name & Official supervisor feedback required. \\
Position & Official supervisor feedback required. \\
Organization & Official supervisor feedback required. \\
Feedback date & Official supervisor feedback required. \\
Evaluation form or source & Official supervisor feedback required. \\
\end{longtable}
\endgroup
}

\vspace{0.5em}\noindent\textbf{Technical feedback}\par

Official supervisor feedback required.

Suggested areas for official technical feedback:

\begin{itemize}
\tightlist
\item
  AWS architecture design
\item
  EKS and Kubernetes deployment
\item
  CI/CD implementation with GitHub Actions and OIDC
\item
  Database, messaging, and event-driven design
\item
  Security and secret handling
\item
  AI service and SageMaker integration
\item
  Monitoring, troubleshooting, and evidence collection
\item
  Final report quality
\end{itemize}

\vspace{0.5em}\noindent\textbf{Working-attitude feedback}\par

Official supervisor feedback required.

Suggested areas for official working-attitude feedback:

\begin{itemize}
\tightlist
\item
  responsibility and ownership
\item
  learning ability
\item
  communication
\item
  teamwork
\item
  discipline and deadline management
\item
  troubleshooting persistence
\item
  documentation habit
\end{itemize}

\vspace{0.5em}\noindent\textbf{Strengths}\par

Official supervisor feedback required.

Based on my own project reflection, the main strengths I demonstrated
were evidence-led troubleshooting, willingness to learn unfamiliar AWS
services, and the ability to connect application code with deployment
and documentation.

\vspace{0.5em}\noindent\textbf{Areas for improvement}\par

Official supervisor feedback required.

Based on my own reflection, I should continue improving IAM policy
design, AWS cost estimation, automated end-to-end testing, and
production evidence collection.

\vspace{0.5em}\noindent\textbf{Overall assessment}\par

Official supervisor feedback required.

No fake names, ratings, quotations, approval statements, or supervisor
conclusions are included in this report.

\vspace{0.5em}\noindent\textbf{Student response}\par

This internship project helped me understand how a modern application
becomes a real AWS system. I learned that successful deployment is not
only about writing application code. It also requires networking, IAM,
Kubernetes readiness, database reliability, asynchronous processing,
logging, monitoring, cost control, and careful documentation.

The most valuable lesson for me was to verify before claiming. A
manifest can show intended configuration, but runtime logs, AWS CLI
outputs, and deployment evidence show what is actually running. I will
continue improving by building stronger test evidence, documenting
operational decisions more clearly, and practicing secure, cost-aware
cloud deployment.
